\documentclass[11pt]{article}

\usepackage[margin=1in]{geometry}
\usepackage[T1]{fontenc}
\usepackage[utf8]{inputenc}
\usepackage{lmodern}
\usepackage{enumitem}
\usepackage{hyperref}
\usepackage{listings}

\hypersetup{colorlinks=true,linkcolor=blue,urlcolor=blue}
\lstset{basicstyle=\ttfamily\small,breaklines=true,columns=fullflexible,keepspaces=true,frame=single}

\title{Outputs, Records, and Visualization}
\author{}
\date{}

\begin{document}
\maketitle

\section{Generation Modes}

The main experiment now has two manual modes:
\begin{itemize}[leftmargin=2em]
    \item \texttt{to\_generate="raw\_data"}: execute the solver, record numerical inspection files, automatically create an independent Results-ready metrics package for each selected configuration, and save only the designated successful frame-by-frame trajectories;
    \item \texttt{to\_generate="visualization"}: read the separately persisted frame-by-frame packages and regenerate Pillow frames without executing MAPF or reading the numerical metrics package.
\end{itemize}

The main experiment no longer has a \texttt{to\_generate="graphs"} mode and does not create PNG metric plots. The reference-comparison workflow remains separate and still supports its own \texttt{raw\_data}, \texttt{graphs}, and \texttt{visualization} modes.

Numerical results and frame-by-frame trajectory data are deliberately separate artifact families. A visualization-only run therefore depends on the corresponding \texttt{frame\_by\_frame} files having already been generated by a raw-data run.

\section{Main-Experiment Outputs}

For every selected exact map configuration, the author-facing numerical inspection file is written under:
\begin{lstlisting}
dev/outputs_main/metrics_data_inspection/
    <numbered_map_category>/
        <numbered_exact_map_configuration>_evaluation.xml
\end{lstlisting}

The program no longer writes a separate \texttt{<map\_config>\_raw\_data.json}.

The same raw-data execution writes an independent Results-ready package for each selected exact map configuration:
\begin{lstlisting}
dev/outputs_main/metrics_data/
    data_dictionary.csv
    <numbered_map_category>/
        <numbered_exact_map_configuration>/
            README.txt
            configuration_metadata.csv
            capacity_summary.csv
            capacity_comparison.csv
            capacity_search_tests.csv
            capacity_search_run_records.csv
            capacity_point_run_records.csv
            capacity_point_summary.csv
            paired_run_comparisons.csv
            results_ready_comparisons.csv
            <numbered_exact_map_configuration>_metrics_data.csv
            metrics_package.json
\end{lstlisting}

The root-level \texttt{data\_dictionary.csv} is the only intentional project-level metrics file. The program does not create cumulative CSVs, a dataset manifest, a root reader guide, or a selected-map-configurations summary file. Running one map configuration updates only that configuration's package. Other configuration folders are neither scanned nor rebuilt. At the beginning of a raw-data run, obsolete root-level cumulative files from the previous design are removed automatically.

The exact-configuration \texttt{\_metrics\_data.csv} is the primary compact table for Results writing. It contains the actual capacity agent numbers, classical and cyclic outcome counts and means at both capacity contexts, absolute and percentage changes, and interpretation flags. The companion files preserve capacity-search decisions, individual run records, paired same-initial-condition comparisons, solver and protocol metadata, variability summaries, seeds, and valid solved-run path counts. No PNG metric plots are produced in this tree.

Terminal output is compartmentalized under:
\begin{lstlisting}
dev/outputs_main/terminal_logs/
    <numbered_map_category>/
        <numbered_exact_map_configuration>/
            raw_data.log
            visualization.log
\end{lstlisting}

\subsection{Selected Frame-by-Frame Runs}

A raw-data execution saves at most two trajectory packages for each of the thirty-two exact configurations:
\begin{itemize}[leftmargin=2em]
    \item the final retained successful classical run at the discovered classical capacity;
    \item the final retained successful cyclic run at the discovered cyclic capacity.
\end{itemize}

Intermediate capacity-search attempts and the cross-mapping comparison runs are not saved as frame-by-frame packages. The directory layout is:
\begin{lstlisting}
dev/outputs_main/frame_by_frame/
    <numbered_map_category>/
        <numbered_exact_map_configuration>/
            shared_context/
                shared_context.pkl
                metadata.json
            classical_capacity_<N>_agents/
                classical/
                    final_selected_successful_run/
                        frame_by_frame.pkl
                        metadata.json
            cyclic_capacity_<N>_agents/
                cyclic/
                    final_selected_successful_run/
                        frame_by_frame.pkl
                        metadata.json
            manifest.json
\end{lstlisting}

The shared context keeps the branch specification and, for dynamic branches, the time-varying map state only once. Each run package keeps the selected agents, solution paths, run configuration, and mapping-specific map information needed by the Pillow renderer. The JSON metadata and manifest make the selection and file locations readable without opening the pickle files.

Recomputing one exact map configuration atomically replaces only that configuration's frame-by-frame subtree. Other map configurations remain intact.

\section{Reference-Comparison Outputs}

Reference numerical data remains under:
\begin{lstlisting}
dev/outputs_ref_comparison/raw_mapf_files/<case_id>/
\end{lstlisting}

The selected trajectories are stored separately by pathfinding mode:
\begin{lstlisting}
dev/outputs_ref_comparison/frame_by_frame/
    single_agent_pf/
        map_1/
            classical/first_successful_run/
            cyclic/first_successful_run/
        map_2/
        map_3/
        manifest.json
    multi_agent_pf/
        map_1/
            classical_capacity_<N>_agents/
                classical/first_successful_run/
                cyclic/first_successful_run/
        map_2/
        map_3/
        manifest.json
\end{lstlisting}

Every terminal run folder contains \texttt{frame\_by\_frame.pkl} and \texttt{metadata.json}. For single-agent pathfinding, the stored package is the first successful timing repetition for each map and mapping. For multi-agent pathfinding, it is the first successful timing repetition from the final repeated comparison at that map's discovered classical capacity. Capacity-search trials are excluded.

Recomputing the single-agent reference case replaces only \texttt{single\_agent\_pf}; recomputing the multi-agent case replaces only \texttt{multi\_agent\_pf}. The other mode remains intact.

\section{Pillow Rendering Contract}

When \texttt{to\_generate="visualization"}, the solver is not called. The renderer loads exactly the packages listed in the relevant frame-by-frame manifest and writes:
\begin{itemize}[leftmargin=2em]
    \item \texttt{config\_000.png}: obstacle-only view;
    \item \texttt{config\_001.png}: showcase view with transitions;
    \item \texttt{config\_002.png}: setup view with agents and targets;
    \item \texttt{frame\_000.png}, \texttt{frame\_001.png}, and so on: the actual solution timeline.
\end{itemize}

For dynamic branches, setup frames use a static setup map derived from the persisted static matrix, while execution frames use the persisted time-varying loop. Campus gray regions remain solver-blocked but may be drawn as visually free cells through the persisted or reconstructed render-only semantic mask.

Main-experiment visualization output mirrors the selected capacity and mapping hierarchy under:
\begin{lstlisting}
dev/outputs_main/visualization/
    <numbered_map_category>/
        <numbered_exact_map_configuration>/
            classical_capacity_<N>_agents/classical/final_selected_successful_run/
            cyclic_capacity_<N>_agents/cyclic/final_selected_successful_run/
            visualization_summary.json
\end{lstlisting}

\section{Failure and Compatibility Behavior}

Each frame-by-frame store has an explicit format version. If the files are missing or incompatible, visualization mode raises an instruction to rerun the relevant case with \texttt{to\_generate="raw\_data"}. Writes use a temporary tree and an atomic subtree replacement so a failed write does not silently destroy the last valid saved copy.

The pickle files are internal project artifacts and should only be loaded from trusted experiment output directories. Their accompanying JSON metadata is intended for inspection, indexing, and troubleshooting.

\end{document}
